% Korean edition of sections/s4_2_oneshare.tex.
\seclabel
\claim{\secnum{4.2}{sec:oneshare}\quad 1주 미만 주문은 0주로 내림되므로 \code{plan()}이 뺀다}

\lead{서버는 주문 금액을 주식 수로 바꿀 때 소수점을 버려 \code{qty2ord}를 정하므로, 1주 비중보다 작은 주문은 0주를 요구하게 된다(\figref{rounding}).} 앞선 주문 4건은 내림과 반올림 어느 쪽에도 맞는다. 삼성전자 0.1\% 주문은 1주 0.0226\%에서 4.4주어치였고 4주를 받았다. 09:15의 0.3\% 주문은 13.3주어치로 13주, LIG넥스원 0.1\% 주문은 1주 0.0426\%에서 2.3주어치로 2주, 08:30의 LIG넥스원 주문은 7.1주어치로 7주였다. 내림 테스트는 두 규칙을 구분한다. 1주가 0.0227\%일 때 삼성전자 0.04\% 주문은 1.76주어치다. 반올림이면 2주, 내림이면 1주다. 서버는 두 번 보냈을 때 모두 \code{qty2ord}를 1로 정했다(주문 ID 1650693, 약 10:26 / 1650719, 10:33:45). 둘 다 즉시 정지해 0주가 체결됐고, 1650719는 주문 직후 \code{Accepted} 상태였다. 서버가 쓴 것으로 보이는 가장 높은 1주 비중 0.0231\%(아래 참조)에서도 이 주문은 1.73주어치이므로 반올림이었다면 여전히 2주였을 것이다. 서버가 시장가보다 20\% 낮은 지정가로 크기를 정했다면 2.2주어치로 어느 규칙이든 2주였을 것이다. 따라서 이 결과는 \code{qty2ord}를 지정가가 아니라 시장가가 정한다는 것도 보여준다.

% Korean edition of figs/rounding.tex.
\fig{rounding}{\code{qty2ord}는 주문 금액을 주식 수로 바꿔 내림한 값이다}[서버가 요구한 주식 수와 반올림 규칙별 예측, 2026년 11월 24일]{%
{\tabfont\begin{tabularx}{\linewidth}{>{\RaggedRight\arraybackslash}X R{15mm} R{9mm} R{12mm} R{12mm} R{14mm} R{12mm}}
\thd{주문} & \thd{1주 (\%)} & \thd{\code{wei} (\%)} & \thd{주식 수 환산} & \thd{내림} & \thd{반올림} & \thd{\code{qty2ord}}\\
LIG넥스원, 08:30 매수(\ref{sec:orders}) & 0.0423 & 0.3 & 7.1 & 7 & 7 & 7\\\rowsep
삼성전자, 09:15 매수(\ref{sec:orders}) & 0.0226 & 0.3 & 13.3 & 13 & 13 & 13\\\rowsep
삼성전자, 매수(\ref{sec:wrappertests}) & 0.0226 & 0.1 & 4.4 & 4 & 4 & 4\\\rowsep
LIG넥스원, 매도(\ref{sec:wrappertests}) & 0.0426 & 0.1 & 2.3 & 2 & 2 & 2\\\rowsep
\hl 삼성전자, 10:26과 10:33 매수, \code{live_rounding}(\ref{sec:oneshare}) & 0.0227 & 0.04 & 1.76 & 1 & 2 & 1\\
\end{tabularx}}}{1주 비중 = 가격 $\times$ 100 / 순자산(\code{bs.na}), 피드 가격 기준, 순자산 대비 \%: 09:52에 0.0226과 0.0426, 10:33에 0.0227. 08:30의 1주 비중은 \ref{sec:orders}절의 7.1주어치에서 구한 0.3 / 7.1이다. 처음 4행은 두 규칙 모두에 맞으므로 둘을 구분할 수 없다. 강조 행이 둘을 구분한다. 1.76주어치는 내림하면 1, 반올림하면 2인데, 서버는 두 주문 모두 1을 요구했다(Id 1650693과 1650719, 즉시 정지, 0주 체결). 09:15 행은 09:52의 1주 비중을 쓴다. 13주는 1주 비중이 0.0214\textasciitilde 0.0231 사이면 모두 맞는다.}{DART Quant Team; \code{tests/live/live_rounding.py}}


\lead{따라서 삼성전자 0.02\% 주문은 HTTP 200으로 접수되면서 0주를 요구하게 된다.} 1주 0.0226\%에서 0.88주어치이므로 0으로 내림되며, 그래도 그날 주문 1건을 쓴다. \code{plan()}은 무엇이든 보내기 전에 이런 주문을 뺀다.

\lead{1주 비중은 가격 × 100 ÷ 순자산이며, 피드가 순자산을 직접 알려준다.} 잔고 레코드 \code{bs}의 \code{na}가 원화 순자산이다. 09:52에 837,002,427원이었다. 모든 델타가 \code{bs_updated}를 싣고 오고 상태가 이를 \code{bs}에 병합하므로, 스냅샷 사이에도 \code{na}는 최신으로 유지된다. 10:33에는 60\,s 안에 844,505,167원에서 844,076,987원으로 움직였다. 서버가 나누는 값이 \code{na}인지는 포지션으로 확인할 수 있다. 보유 종목마다 \code{pos} $\times$ \code{close} $\times$ 100 / \code{wei}로 추정값을 얻는데, 09:52 보유 종목 23개의 중앙값은 836,798,658원으로 \code{na}보다 0.02\% 낮았다. \code{wei}가 소수 둘째 자리까지이므로 추정값 쪽이 오차가 크다. 0.43\%인 LIG넥스원은 반올림 오차 0.005만으로도 추정값이 1.2\% 움직인다. 그래서 \code{nav()}는 순자산 \code{bs.na}를 돌려주며, 피드의 별도 필드 \code{nav}는 쓰지 않는다. \code{nav()}는 피드에 \code{na}가 없을 때만 중앙값 추정으로 대체한다(D24). 09:52에 1주 비중은 삼성전자 0.0226\%, LIG넥스원 0.0426\%, SK하이닉스 0.1413\%였다.

\lead{보유 종목의 가격은 호출자가 주지 않으면 피드의 \code{close}를 쓰고, 그 밖의 종목은 호출자가 가격을 줘야 한다.} 대회 서버는 포트폴리오 밖 종목의 가격을 알려주지 않고, 시장 데이터는 래퍼 밖에 둔다(N1). 그래서 \code{plan()}은 종목코드에서 가격으로 가는 딕셔너리 \code{prices=}를 받는다. 스냅샷 필드 \code{noPosPrcs}가 포트폴리오 밖 종목의 가격인지는 확인하지 않았으며, 확인할 때까지는 호출자가 \code{prices=}를 넘긴다. 가격이 없으면 확인하지 못했다는 경고를 남기고 주문을 유지한다. 래퍼가 잴 수 없는 주문을 빼면 근거 없이 전략의 목표를 바꾸게 되기 때문이다. 뺀 주문은 하나씩 로그에 남기고 \code{tf.dropped}에 보관하므로 호출자가 무엇이 빠졌는지 볼 수 있다(D22).

\lead{09:52에 필터는 1주 비중보다 작은 주문 1건을 빼고, 큰 주문 2건은 남겼다.} 목표는 현재 비중에 SK하이닉스 0.03\%p, 삼성전자 0.1\%p, 리노공업 0.06\%p를 더한 것이었다. 0.03 주문이 너무 작다고 건너뛰어지지 않고 필터까지 오도록 \code{min_wei}를 기본값 0.05에서 0.01로 낮췄다. SK하이닉스 +0.03\%는 1주가 0.1413\%여서 경고와 함께 빠졌고, 삼성전자 +0.1\%(1주 0.0226\%)와 리노공업 +0.06\%는 남았다.

\lead{필터는 여유를 두지 않는다.} 1주 비중보다 몇 \% 큰 주문도 \code{plan()}과 서버가 \code{qty2ord}를 계산하는 순간 사이에 가격이 오르면 0으로 내림될 수 있다. \figref{rounding}의 09:15 주문에서 서버가 쓴 1주 비중은 0.0214\%와 0.0231\% 사이로, 8\%의 범위였다. 고정 여유를 두면 서버가 1주로 정할 주문까지 빠지므로, 래퍼는 이 선택을 호출자에게 맡긴다. 호출자는 \code{min_wei}를 올릴 수 있다.

\verdict{래퍼의 \code{plan()}은 1주 비중을 가격 $\times$ 100 / \code{bs.na}로 계산한다. 가격은 \code{prices=}가 있으면 그 값, 없으면 보유 종목에 한해 피드의 \code{close}다. 그보다 작은 주문은 경고와 함께 빼서 \code{tf.dropped}에 보관하고, 가격을 알 수 없는 주문은 경고와 함께 남긴다(D20, D22, D24).}
